\documentclass[10pt]{article}
\usepackage[a4paper,margin=1.6cm]{geometry}
\usepackage{amsmath,amssymb,booktabs,enumitem}
\usepackage[most]{tcolorbox}
\setlist{nosep,leftmargin=1.2em}
\setlength{\parindent}{0pt}\setlength{\parskip}{3pt}
\pagestyle{empty}
\begin{document}

{\large\bfseries The Route Law of Multimodal Forgetting}\\[1pt]
{\small Two modalities, class- and domain-incremental learning. Working note.}

\textbf{Setup.} A model with two modalities learns task~1, then keeps training on task~2. Each task-1 sample is typed by
its pointwise information decomposition (PID; $I_{\text{ccs}}$ redundancy): \emph{redundant} (R: either modality alone
gives the label), \emph{unique} (U$_1$, U$_2$: only that modality does), or \emph{synergistic} (S: only the two
together do). For each type $\tau$, the \emph{retention} $\rho_\tau$ is the share of correctly learned task-1 test samples
still correct after task~2 (chance-corrected).

\textbf{Idea.} Each modality is a \emph{route} that carries old knowledge through the new task. Let $p_m$ be the
probability that route $m$ survives. Unique knowledge needs its own route, redundant knowledge needs \emph{any} route,
synergistic knowledge needs \emph{all} routes.

\begin{tcolorbox}[colback=gray!6,colframe=black!55,boxrule=0.5pt,left=4pt,right=4pt,top=2pt,bottom=2pt]
\textbf{Route law} (independent route failure; $p_m$ measured as $\rho_{\text{U}_m}$, no fitted parameters)
\[
\rho_{\text{U}_m} = p_m, \qquad
\rho_{\text{R}} = 1-(1-p_1)(1-p_2), \qquad
\rho_{\text{S}} = p_1\,p_2 .
\]
\textbf{With a shared bottleneck} (the fusion layer, survival $h$): every retention is multiplied by $h$.
\end{tcolorbox}

\textbf{Consequences.}
\begin{enumerate}[label=\arabic*.]
\item \emph{Order:} $\rho_{\text{S}} \le \rho_{\text{U}_m} \le \rho_{\text{R}}$. Synergy is forgotten first, redundancy last.
\item \emph{Size of the gap:} $\rho_{\text{R}}-\rho_{\text{S}} = p_1+p_2-2p_1p_2$ (for $p_1{=}p_2{=}p$: $2p(1-p)$). The gap vanishes
when routes rarely fail ($p\!\to\!1$) and is largest at $p=\tfrac12$. So the effect depends on how hard the new task hits
each route, not on the dataset alone.
\item \emph{Locality:} if task~2 changes only modality~$m$, then only U$_m$ and S drop; R survives on the other route.
\item \emph{Identity:} $\rho_{\text{R}}+\rho_{\text{S}}=\rho_{\text{U}_1}+\rho_{\text{U}_2}$, for any shared bottleneck $h$; and $h=\rho_{\text{U}_1}\rho_{\text{U}_2}/\rho_{\text{S}}$ becomes measurable.
\end{enumerate}

\textbf{Evidence.} Synthetic data with known types (5 robustness checks: equal modality strength, equal type mix,
single-modality shifts, both combined; 5 seeds) and frozen-encoder features of three real datasets (two-view CIFAR-100,
AV-MNIST, CREMA-D; 2 class splits, 2 training recipes, 3 seeds); with and without replay.

\begin{minipage}[t]{0.49\linewidth}
\small\centering
\textbf{(a)} Zero-parameter prediction vs.\ observation\\[2pt]
\begin{tabular}{lccc}
\toprule
data & runs & R: $r$ / $|$err$|$ & S: $r$ / $|$err$|$\\
\midrule
synthetic & 57 & 0.93 / 0.08 & 0.94 / 0.06\\
real features & 60 & 0.94 / 0.07 & 0.89 / 0.12\\
\bottomrule
\end{tabular}\\[5pt]
\textbf{(c)} Real data, with replay (observed / predicted)\\[2pt]
\begin{tabular}{lccc}
\toprule
dataset & $p_1,p_2$ & $\rho_{\text{R}}$ & $\rho_{\text{S}}$\\
\midrule
two-view CIFAR & .57, .56 & .79 / .81 & .32 / .32\\
CREMA-D & .87, .90 & .96 / .99 & .77 / .79\\
\bottomrule
\end{tabular}
\end{minipage}\hfill
\begin{minipage}[t]{0.49\linewidth}
\small\centering
\textbf{(b)} Locality, synthetic, no replay (observed / predicted)\\[2pt]
\begin{tabular}{lcc}
\toprule
 & only $m_1$ changes & only $m_2$ changes\\
\midrule
U$_1$ & .26 (lost) & 1.00 (kept)\\
U$_2$ & .96 (kept) & .21 (lost)\\
R & .57 / .97 & .99 / 1.00\\
S & .27 / .25 & .19 / .21\\
\bottomrule
\end{tabular}\\[5pt]
\textbf{(d)} Difficulty control, synthetic, replay\\[2pt]
\begin{tabular}{lc}
\toprule
check & S vs.\ R at equal confidence\\
\midrule
original / equal type mix & 0.50 / 0.49\\
equal strength / both & 0.51 / 0.49\\
\bottomrule
\end{tabular}
\end{minipage}

\vspace{3pt}
\begin{itemize}
\item Synergy is lost about \textbf{twice as fast} as redundancy even at equal starting confidence (d), so the effect is about
the \emph{type} of knowledge, not how well it was learned.
\item The law explains why the effect looked absent on CREMA-D: routes rarely fail there ($p\approx0.9$), so the gap is small (2); on
two-view CIFAR routes fail often ($p\approx0.57$) and the gap is large (c).
\item Measured bottleneck: $h\approx0.9$--$1.0$ with replay, $h\approx0.2$--$0.3$ without. Replay mainly protects the shared fusion layer.
\end{itemize}

\textbf{Where it fails / limits.} Without replay, R is over-predicted when one modality dominates (b, left): the model stores
redundant knowledge mostly in the strong route, so the backup route is weak. Real-data synergy is rare (2--8\% of samples) and its
labels come from probes (36\% agreement with human ratings on CREMA-D). Synthetic synergy is a modular sum. Frozen encoders; two tasks.

\textbf{Next.} Three modalities (law: $\rho_{\text{S}_{123}}=p_1p_2p_3$, $\rho_{\text{R}_{123}}=1-\prod_m(1-p_m)$); three-way synergy is learnable
with 3 classes. Interventions that change one route on real data. Weighting routes by how much redundant knowledge each actually stores.

\end{document}
